% Source: https://artofproblemsolving.com/community/c3h1760770_chinese_remainder_theorem
% Stable identity: aops-topic-1760770-q01
\textbf{Statement.} If $m$ and $n$ are relatively prime positive integers, then for any integers $a,b$ the system $x\equiv a\pmod m$, $x\equiv b\pmod n$ has a solution, and the solutions form exactly one residue class modulo $mn$. (Reason: the $n$ numbers $a,a+m,\dots,a+(n-1)m$ have different remainders modulo $n$, because $n\mid(i-j)m$ forces $n\mid i-j$; so exactly one of them is $\equiv b\pmod n$.) Applying it repeatedly handles any number of pairwise relatively prime moduli.

\textbf{The example.} $x\equiv2\pmod3$ and $x\equiv2\pmod7$ together say $x\equiv2\pmod{21}$, so $x\in\{2,23,44,65,86,\dots\}$. The first of these that is $\equiv3\pmod5$ is $23$. All solutions are $x\equiv23\pmod{105}$, and the smallest positive one is $\boxed{23}$.
